\documentclass[10pt,a4paper]{amsart}
\usepackage[margin=19mm]{geometry}
\usepackage{amsmath,amssymb,mathtools}
\usepackage[scr=rsfs,cal=euler]{mathalfa}
\usepackage{xfrac}
\usepackage[unicode,hidelinks,bookmarks=false]{hyperref}
\newcommand{\ie}{i.e.\ }
\newcommand{\cf}{cf.\ }
\newcommand{\resp}{respectively\ }
\newcommand{\Subsection}{\subsection}
\providecommand{\nobreakdash}{\nobreak\hbox{-}}
\setlength{\emergencystretch}{2em}
\multlinegap0pt
\usepackage{amssymb}
\usepackage{enumerate}
\usepackage{xcolor}
\usepackage{tikz}
\newtheorem{lemma}{Lemma}
\title{Existence of weak solutions for nonlinear drift-diffusion equations with measure data}
\author{Sukjung Hwang, Kyungkeun Kang, Hwa Kil Kim, Jung-Tae Park}
\date{Source: arXiv:2501.07847v2 (CC BY 4.0)}
\begin{document}
\maketitle

\noindent\emph{Source and licence.} Existence of weak solutions for nonlinear drift-diffusion equations with measure data. Version arXiv:2501.07847v2 (CC BY 4.0), distributed by arXiv under the Creative Commons Attribution 4.0 International licence, http://creativecommons.org/licenses/by/4.0/, as recorded in the arXiv licence metadata for this exact version and shown on the abstract page https://arxiv.org/abs/2501.07847v2. Full licence notice: \texttt{licenses/arxiv-2501.07847-CC-BY-4.0.txt}.\par\smallskip

\noindent\textit{Editorial scope.} Counted unit: the article's lemma L1 (source0.tex lines 1044--1057). The article's complete original proof of it is delivered with it (source0.tex lines 1059--1640, one \texttt{proof} environment of 582 lines). No mathematics is written, repaired or re-derived: the statement and the proof are the article's own lines, in the article's own notation, and no retained line is substituted or re-flowed.  Retained uncounted as the source-local context the proof needs: lines 292--294.  Nothing was omitted from any retained span, so the package carries no omission record.  No retained line contains a citation, so the wrapper carries no bibliography and no bibliography entry.  No retained line contains a figure environment, an image-inclusion command or drawing code, so no image is delivered and the package declares no asset dependency.  Editorial formatting only, in addition: an AMS \texttt{amsart} preamble that restates the article's own declarations and macros used by the retained lines, namely the environments \texttt{lemma} on separate counters, together with the standard packages those lines need.  The retained environments print in this document as Lemma 1 in order of appearance, whereas the article prints the same items with the numbering of its own full document.  The complete native source of the article as distributed in the arXiv e-print for this exact version is bundled unchanged as \texttt{sources/171581/source0.tex}.
\begin{equation}\label{PME}
	u_t - \Delta u^m + \nabla \cdot \left( u \, V \right)= \mu, \ \text{ for } \  m>0 \ \text{ on } \ \Omega_{T},
\end{equation}

\begin{lemma}\label{L1}
Let $d\geq 2$ and $1-\frac{1}{d} < m \leq 2$.
Suppose that $u$ is a regular solution of \eqref{PME}.
Then it follows that
\begin{equation}\label{Estimate01}
	\sup_{t\in (0, T)} \int_{\Omega} u(x, t) \,dx \leq \mu (\Omega_{T}),
\end{equation}
where $\mu(\Omega_T): = \int_{\Omega_T} \mu(x,t) \, dxdt$. 
Moreover, for any $A>0$, $\xi > 1$, and $1 < q \le m+1$, we have
\begin{equation}\label{Estimate02}
	\int_{\Omega_T} \frac{\left|\nabla u^{\frac{m+q-1}{2}}\right|^2}{\left(A^{q-1} + u^{q-1}\right)^{\xi}} \,dxdt		
	\leq \frac{(m+q-1)^2 A^{(q-1)(1-\xi)}}{m(q-1)(\xi-1)} \left[ 2\mu(\Omega_{T}) + \frac{(q-1)(\xi-1)}{3m}\int_{\Omega_T} u^{2-m} |V|^2 \,dxdt\right].
\end{equation}
\end{lemma}

\begin{proof}
We first note that, since $d\ge2$ and $m>1-\frac1d$, we have
$2-m\le m+1$. Hence, from
$u\in C([0,T];L^{m+1}(\Omega))$ and $V\in L^\infty(\Omega_T)$, it follows that $u^{2-m}|V|^2\in L^1(\Omega_T)$.

We justify the nonlinear test functions by a standard smooth approximation argument. More precisely, we first perform the following computations for smooth non-degenerate approximating solutions $u_\kappa$ with zero lateral boundary data and zero initial data. For such solutions, all chain rules and integrations by parts below are classical, and the nonlinear test functions are admissible. The estimates obtained are independent of the approximation parameter $\kappa$. Hence, passing to the limit $\kappa\to0$, using the strong convergence of the approximations and the lower semicontinuity of the quadratic gradient terms, yields the estimates for the original regular solution. To simplify notation, we omit the subscript $\kappa$ below.


%For the smooth approximating solutions, the weak formulation is equivalently
%written as
%\[
%\int_{\Omega_T}
%\left\{
%u_t\eta+\nabla u^m\cdot\nabla\eta-uV\cdot\nabla\eta
%\right\}\,dxdt
%=
%\int_{\Omega_T}\mu\eta\,dxdt
%\]
%for admissible test functions $\eta$ vanishing on
%$\partial\Omega\times(0,T)$ and satisfying $\eta(\cdot,T)=0$.

%The restriction $q\le m+1$ is used to justify the nonlinear powers appearing in the test functions from the assumption $u^m\in L^2(0,T;W^{1,2}_0(\Omega))$.
%Indeed, on the truncated sets where $u>\delta$, the identities
%\[
%\nabla u^{q-1} = \frac{q-1}{m}u^{q-1-m}\nabla u^m
%\]
%and
%\[
%\nabla u^{\frac{m+q-1}{2}} = \frac{m+q-1}{2m} u^{\frac{q-m-1}{2}}\nabla u^m
%\]
%are controlled by $\nabla u^m$, since $q-1-m\le0$ and $(q-m-1)/2\le0$.



\smallskip
\noindent
\underline{\emph{Step 1: proof of \eqref{Estimate01}.}}

Fix $\varepsilon>0$ and $0<\delta<\varepsilon$. Define
\[
F_{\varepsilon,\delta}(s):=
\begin{cases}
0, & 0\le s\le \delta,\\[1mm]
\dfrac{s^{q-1}-\delta^{q-1}}
{\varepsilon^{q-1}-\delta^{q-1}},
& \delta<s<\varepsilon,\\[3mm]
1, & s\ge \varepsilon,
\end{cases} 
\quad \text{ and }
\quad 
\Psi_{\varepsilon,\delta}(s)
:=
\int_0^sF_{\varepsilon,\delta}(r)\,dr.
\]
Let $\phi\in C^\infty([0,T])$ satisfy $0\le \phi\le1$, $-\phi_t\ge0$, $\phi(T)=0$.
We use
\[
\eta_1:=F_{\varepsilon,\delta}(u)\phi
\]
as a test function. Since $u=0$ on $\partial\Omega\times(0,T)$ and
$F_{\varepsilon,\delta}(0)=0$, we have $\eta_1=0$ on $\partial\Omega\times(0,T)$.
Thus $\eta_1$ is admissible.

Testing the equation by $\eta_1$, we obtain
\[
I+II=III+IV,
\]
where
\[
I:=
\int_{\Omega_T}
u_tF_{\varepsilon,\delta}(u)\phi\,dxdt,
\qquad
II:=
\int_{\Omega_T}
\nabla u^m\cdot\nabla F_{\varepsilon,\delta}(u)\phi\,dxdt,
\]
\[
III:=
\int_{\Omega_T}
\mu F_{\varepsilon,\delta}(u)\phi\,dxdt,
\qquad
IV:=
\int_{\Omega_T}
uV\cdot\nabla F_{\varepsilon,\delta}(u)\phi\,dxdt.
\]

Since $\partial_t\Psi_{\varepsilon,\delta}(u) = u_tF_{\varepsilon,\delta}(u)$, 
integrating by parts in time gives
\[
I
=
\int_{\Omega_T}
\Psi_{\varepsilon,\delta}(u)(-\phi_t)\,dxdt 
-\int_\Omega
\Psi_{\varepsilon,\delta}(u(x,0))\phi(0)\,dx
+
\int_\Omega
\Psi_{\varepsilon,\delta}(u(x,T))\phi(T)\,dx.
\]
The terminal term is zero because $\phi(T)=0$. The initial term is also zero,
because $u(\cdot,0)=0$ and 
$\Psi_{\varepsilon,\delta}(0)=0$.
Therefore
\begin{equation}\label{I01-smooth-u}
I
=
\int_{\Omega_T}
\Psi_{\varepsilon,\delta}(u)(-\phi_t)\,dxdt
\ge0.
\end{equation}

Moreover, since $0\le F_{\varepsilon,\delta}\le1$ and $0\le\phi\le1$,
\begin{equation}\label{III01-smooth-u}
0\le III\le \mu(\Omega_T).
\end{equation}

Next,
\[
\nabla F_{\varepsilon,\delta}(u)
=
\frac{(q-1)u^{q-2}}
{\varepsilon^{q-1}-\delta^{q-1}}
\chi_{\{\delta<u<\varepsilon\}}
\nabla u.
\]
Hence
\[
\begin{aligned}
II
&=
\frac{m(q-1)}
{\varepsilon^{q-1}-\delta^{q-1}}
\int_{\{\delta<u<\varepsilon\}}
u^{m+q-3}
|\nabla u|^2\phi\,dxdt  \\
&=
\frac{4m(q-1)}
{(m+q-1)^2(\varepsilon^{q-1}-\delta^{q-1})}
\int_{\{\delta<u<\varepsilon\}}
\left|\nabla u^{\frac{m+q-1}{2}}\right|^2\phi\,dxdt.
\end{aligned}
\]
In particular, $II\ge0$. Letting $\delta\to0$ and using Fatou's lemma, we get
\begin{equation}\label{II01-smooth-u}
\liminf_{\delta\to0}II
\ge
\frac{4m(q-1)}
{(m+q-1)^2\varepsilon^{q-1}}
\int_{\{0<u<\varepsilon\}}
\left|\nabla u^{\frac{m+q-1}{2}}\right|^2\phi\,dxdt.
\end{equation}

For the drift term, Young's inequality gives
\[
\begin{aligned}
IV
&=
\frac{q-1}{\varepsilon^{q-1}-\delta^{q-1}}
\int_{\{\delta<u<\varepsilon\}}
u^{q-1}V\cdot\nabla u\,\phi\,dxdt  \\
&\le
\frac12 II
+
\frac{q-1}{2m(\varepsilon^{q-1}-\delta^{q-1})}
\int_{\{\delta<u<\varepsilon\}}
u^{q+1-m}|V|^2\phi\,dxdt.
\end{aligned}
\]
On $\{\delta<u<\varepsilon\}$, we have
$
u^{q+1-m}
=
u^{2-m}u^{q-1}
\le
\varepsilon^{q-1}u^{2-m}$.
Therefore
\begin{equation}\label{IV01-smooth-u}
IV
\le
\frac12 II
+
\frac{(q-1)\varepsilon^{q-1}}
{2m(\varepsilon^{q-1}-\delta^{q-1})}
\int_{\{\delta<u<\varepsilon\}}
u^{2-m}|V|^2\phi\,dxdt.
\end{equation}

Combining \eqref{I01-smooth-u}--\eqref{IV01-smooth-u}, and then letting
$\delta\to0$, we obtain
\begin{equation}\label{comp1-smooth-u}
\begin{aligned}
&\int_{\Omega_T}
\left[
\int_0^u
\min\left\{1,\frac{s^{q-1}}{\varepsilon^{q-1}}\right\}\,ds
\right]
(-\phi_t)\,dxdt
+
\frac{2m(q-1)}
{(m+q-1)^2\varepsilon^{q-1}}
\int_{\{0<u<\varepsilon\}}
\left|\nabla u^{\frac{m+q-1}{2}}\right|^2\phi\,dxdt  \\
&\qquad\le
\mu(\Omega_T)
+
\frac{q-1}{2m}
\int_{\{0<u<\varepsilon\}}
u^{2-m}|V|^2\phi\,dxdt.
\end{aligned}
\end{equation}

Letting $\varepsilon\to0$, the last term on the right-hand side tends to
zero because $u^{2-m}|V|^2\in L^1(\Omega_T)$. Also,
\[
\int_0^u
\min\left\{1,\frac{s^{q-1}}{\varepsilon^{q-1}}\right\}\,ds
\longrightarrow u
\qquad\text{in }L^1(\Omega_T).
\]
Thus
\[
\int_{\Omega_T}u(-\phi_t)\,dxdt
\le
\mu(\Omega_T)
\]
for every smooth nonincreasing $\phi$ satisfying
$0\le\phi\le1$ and $\phi(T)=0$.

Now fix $\tau\in(0,T)$. Choose nonnegative functions
$\rho_j\in C_c^\infty(0,T)$ such that
\[
\int_0^T\rho_j(t)\,dt=1,
\qquad
\rho_j\rightharpoonup\delta_\tau
\]
in the sense of measures, and set
\[
\phi_j(t):=\int_t^T\rho_j(s)\,ds.
\]
Then $0\le\phi_j\le1$, $-(\phi_j)_t=\rho_j$, and $\phi_j(T)=0$.
Since $u\in C([0,T];L^{m+1}(\Omega))$, the function
\[
t\mapsto\int_\Omega u(x,t)\,dx
\]
is continuous. Passing to the limit $j\to\infty$ gives
\[
\int_\Omega u(x,\tau)\,dx
\le
\mu(\Omega_T).
\]
Taking the supremum over $\tau\in(0,T)$ proves \eqref{Estimate01}.

\bigskip
\noindent
\underline{\emph{Step 2: proof of \eqref{Estimate02}.}}

Let $A>0$ and $\xi>1$. We again fix $\varepsilon>0$ and
$0<\delta<\varepsilon$, and use the test function
\[
\eta_2:=
\frac{F_{\varepsilon,\delta}(u)\phi}
{\left(A^{q-1}+u^{q-1}\right)^{\xi-1}},
\]
where $\phi\in C^\infty([0,T])$ satisfies $0\le\phi\le1$, $ -\phi_t\ge0$, and $\phi(T)=0$.
Since $F_{\varepsilon,\delta}(0)=0$, the function $\eta_2$ vanishes on
$\partial\Omega\times(0,T)$, and hence it is admissible.

Testing the equation by $\eta_2$ gives
\[
I+II_1+II_2=III+IV_1+IV_2,
\]
where
\[
I:=
\int_{\Omega_T}
u_t
\frac{F_{\varepsilon,\delta}(u)\phi}
{\left(A^{q-1}+u^{q-1}\right)^{\xi-1}}
\,dxdt,
\qquad
II_1:=
\int_{\Omega_T}
\frac{
\nabla u^m\cdot\nabla F_{\varepsilon,\delta}(u)
}
{\left(A^{q-1}+u^{q-1}\right)^{\xi-1}}
\phi\,dxdt,
\]
\[
II_2:=
(1-\xi)
\int_{\Omega_T}
\frac{
F_{\varepsilon,\delta}(u)
\nabla u^m\cdot\nabla u^{q-1}
}
{\left(A^{q-1}+u^{q-1}\right)^\xi}
\phi\,dxdt,
\qquad
III:=
\int_{\Omega_T}
\mu
\frac{F_{\varepsilon,\delta}(u)\phi}
{\left(A^{q-1}+u^{q-1}\right)^{\xi-1}}
\,dxdt,
\]
\[
IV_1:=
\int_{\Omega_T}
\frac{
uV\cdot\nabla F_{\varepsilon,\delta}(u)
}
{\left(A^{q-1}+u^{q-1}\right)^{\xi-1}}
\phi\,dxdt,
\qquad
IV_2:=
(1-\xi)
\int_{\Omega_T}
\frac{
F_{\varepsilon,\delta}(u)
uV\cdot\nabla u^{q-1}
}
{\left(A^{q-1}+u^{q-1}\right)^\xi}
\phi\,dxdt.
\]

Since $F_{\varepsilon,\delta}(u)=0$ on $\{u\le\delta\}$, all terms
involving $\nabla u^{q-1}$ are supported in $\{u>\delta\}$. Hence no
singularity occurs at $u=0$. Moreover,
\[
\nabla u^m\cdot\nabla u^{q-1}
=
m(q-1)u^{m+q-3}|\nabla u|^2
\ge0
\qquad\text{on }\{u>\delta\}.
\]
Since $1-\xi<0$, we have
\[
II_2\le0.
\]
Therefore
\begin{equation}\label{basic-smooth-u}
-II_2
\le
|I|+II_1+III+|IV_1|+|IV_2|.
\end{equation}

Define
\[
\Theta_{\varepsilon,\delta}(s)
:=
\int_0^s
\frac{F_{\varepsilon,\delta}(r)}
{\left(A^{q-1}+r^{q-1}\right)^{\xi-1}}
\,dr.
\]
Then
\[
I
=
\int_{\Omega_T}
\partial_t\Theta_{\varepsilon,\delta}(u)\phi\,dxdt.
\]
As in Step~1, the initial and terminal boundary terms vanish, and hence
\[
I
=
\int_{\Omega_T}
\Theta_{\varepsilon,\delta}(u)(-\phi_t)\,dxdt.
\]
Moreover,
\[
0\le
\Theta_{\varepsilon,\delta}(u)
\le
A^{(q-1)(1-\xi)}u.
\]
Using \eqref{Estimate01}, we get
\begin{equation}\label{I2-smooth-u}
|I|
\le
A^{(q-1)(1-\xi)}\mu(\Omega_T).
\end{equation}
Similarly,
\[
0\le
\frac{F_{\varepsilon,\delta}(u)}
{\left(A^{q-1}+u^{q-1}\right)^{\xi-1}}
\le
A^{(q-1)(1-\xi)},
\]
so
\begin{equation}\label{III2-smooth-u}
0\le III
\le
A^{(q-1)(1-\xi)}\mu(\Omega_T).
\end{equation}

For $II_1$, we have
\[
\begin{aligned}
II_1
&=
\frac{m(q-1)}
{\varepsilon^{q-1}-\delta^{q-1}}
\int_{\{\delta<u<\varepsilon\}}
\frac{
u^{m+q-3}|\nabla u|^2\phi
}
{\left(A^{q-1}+u^{q-1}\right)^{\xi-1}}
\,dxdt \\
&\le
A^{(q-1)(1-\xi)}
\frac{m(q-1)}
{\varepsilon^{q-1}-\delta^{q-1}}
\int_{\{\delta<u<\varepsilon\}}
u^{m+q-3}|\nabla u|^2\phi
\,dxdt.
\end{aligned}
\]
The last integral is precisely the diffusion term from Step~1. Hence, using
\eqref{comp1-smooth-u} before letting $\varepsilon\to0$, we obtain
\begin{equation}\label{II1-smooth-u}
\begin{aligned}
\limsup_{\delta\to0}II_1
&\le
2A^{(q-1)(1-\xi)}
\left[
\mu(\Omega_T)
+
\frac{q-1}{2m}
\int_{\{0<u<\varepsilon\}}
u^{2-m}|V|^2\phi\,dxdt
\right].
\end{aligned}
\end{equation}

Next, by Fatou's lemma,
\begin{equation}\label{II2-smooth-u}
\liminf_{\varepsilon\to0}
\liminf_{\delta\to0}
(-II_2) \ge
\frac{4m(q-1)(\xi-1)}
{(m+q-1)^2}
\int_{\Omega_T}
\frac{
\left|\nabla u^{\frac{m+q-1}{2}}\right|^2\phi
}
{\left(A^{q-1}+u^{q-1}\right)^\xi}
\,dxdt.
\end{equation}

For $IV_1$, Young's inequality gives
\[
\begin{aligned}
|IV_1|
&=
\left|
\frac{q-1}{\varepsilon^{q-1}-\delta^{q-1}}
\int_{\{\delta<u<\varepsilon\}}
\frac{
u^{q-1}V\cdot\nabla u\,\phi
}
{\left(A^{q-1}+u^{q-1}\right)^{\xi-1}}
\,dxdt
\right| \\
&\le
II_1 +
\frac{q-1}{4m(\varepsilon^{q-1}-\delta^{q-1})}
\int_{\{\delta<u<\varepsilon\}}
\frac{
u^{q+1-m}|V|^2\phi
}
{\left(A^{q-1}+u^{q-1}\right)^{\xi-1}}
\,dxdt.
\end{aligned}
\]
Since
$
\left(A^{q-1}+u^{q-1}\right)^{1-\xi}
\le
A^{(q-1)(1-\xi)}
$
and
$
u^{q+1-m}
=
u^{2-m}u^{q-1}
\le
\varepsilon^{q-1}u^{2-m}
$ 
on $\{\delta<u<\varepsilon\}$,
we infer
\begin{equation}\label{IV1-smooth-u}
\limsup_{\delta\to0}|IV_1|
\le
\limsup_{\delta\to0}II_1
+
o_\varepsilon(1),
\end{equation}
where $o_\varepsilon(1)\to0$ as $\varepsilon\to0$.

For $IV_2$, Young's inequality yields
\[
|IV_2|
\le
\frac14(-II_2)
+
\frac{(q-1)(\xi-1)}{m}
\int_{\Omega_T}
\frac{
F_{\varepsilon,\delta}(u)u^{q+1-m}|V|^2\phi
}
{\left(A^{q-1}+u^{q-1}\right)^\xi}
\,dxdt.
\]
Using
$
\left(A^{q-1}+u^{q-1}\right)^\xi
\ge
A^{(q-1)(\xi-1)}u^{q-1}
$ on $\{u>0\}$,
we obtain
\[
\frac{u^{q+1-m}}
{\left(A^{q-1}+u^{q-1}\right)^\xi}
\le
A^{(q-1)(1-\xi)}u^{2-m}.
\]
Therefore
\begin{equation}\label{IV2-smooth-u}
|IV_2|
\le
\frac14(-II_2)
+
\frac{(q-1)(\xi-1)A^{(q-1)(1-\xi)}}{m}
\int_{\Omega_T}u^{2-m}|V|^2\phi\,dxdt.
\end{equation}

Combining
\eqref{basic-smooth-u}--\eqref{IV2-smooth-u}, then letting
$\delta\to0$ and $\varepsilon\to0$, and using
\[
\int_{\{0<u<\varepsilon\}}
u^{2-m}|V|^2\phi\,dxdt
\to0
\qquad\text{as }\varepsilon\to0,
\]
we get
\[
\int_{\Omega_T}
\frac{
\left|\nabla u^{\frac{m+q-1}{2}}\right|^2\phi
}
{\left(A^{q-1}+u^{q-1}\right)^\xi}
\,dxdt \le
\frac{(m+q-1)^2A^{(q-1)(1-\xi)}}
{m(q-1)(\xi-1)}
\left[
2\mu(\Omega_T)
+
\frac{(q-1)(\xi-1)}{3m}
\int_{\Omega_T}u^{2-m}|V|^2\phi\,dxdt
\right].
\]


Finally, choose $\{\phi_k\}_{k=1}^\infty\subset C^\infty([0,T])$ such that
\[
0\le\phi_k\le1,\qquad -(\phi_k)_t\ge0,\qquad \phi_k(T)=0,
\]
\[
\phi_k(t)=1
\quad\text{for }0\le t\le T-\frac1k,
\qquad
\phi_k\nearrow1
\quad\text{pointwise in }(0,T).
\]
Applying the previous estimate with $\phi=\phi_k$ and then letting
$k\to\infty$, the monotone convergence theorem gives \eqref{Estimate02}, which completes the proof.
\end{proof}

\end{document}
